\section{KoalaFalcon}

\subsection{Core Construction} 

KoalaFalcon is a variant of CoreFalcon+ (itself a variant of Falcon 
with a slight modification to better enable security analysis with Rényi divergence). 
The primary algorithmic difference in KoalaFalcon is to use Poseidon as the hash function 
instead of SHAKE. Beyond this, KoalaFalcon fixes the modulus $q$ to the KoalaBear prime. 
To support the KoalaBear prime, we set the following formulaic parameters from CoreFalcon+, 
explaining the rationale and applicability of this choice in \hyperref[sqrtq-scaling-argument]{\ref*{sqrtq-scaling-argument}}:
\begin{align*}
s &= \frac{1}{\pi} \sqrt{\frac{\ln\!\left(4n\left(1+\frac{1}{\epsilon}\right)\right)}{2}} \cdot \alpha \sqrt{q}, &
\beta &= \tau s \sqrt{2n}.
\end{align*}

% TODO cite corefalcon+

KoalaFalcon's signing and verification algorithms are depicted in
\hyperref[fig:koalafalcon-construction]{Figure~\ref*{fig:koalafalcon-construction}}, 
with two concrete parameter sets defined for KoalaFalcon in 
\hyperref[tab:koalafalcon-parameters]{Table~\ref*{tab:koalafalcon-parameters}}.
Algorithmically, the primary difference between CoreFalcon+ and KoalaFalcon is 
in the hash function. KoalaFalcon is able to use Poseidon natively because it 
fixes $q$ as the KoalaBear prime, supporting pre-existing Poseidon instantiations 
under the same modulus. Messages are hashed into $R_q$ elements by taking repeated Poseidon 
permutations, collecting the sponge state between each invocation until $n$ 
coefficients have been collected. The rest remains algorithmically unchanged in 
KoalaFalcon aside from re-parameterization.

\begin{figure}[!htbp] \label{koalafalcon-construction}
  \centering
  \setlength{\fboxsep}{7pt}
  \fbox{%
  \begin{minipage}{0.94\textwidth}
  \renewcommand{\arraystretch}{0.95}
  \begin{tabular}{@{\hspace{0.5em}}>{\raggedright\arraybackslash}p{\dimexpr0.5\linewidth-0.85em-0.5\arrayrulewidth\relax}@{\hspace{0.35em}}|@{\hspace{0.35em}}>{\raggedright\arraybackslash}p{\dimexpr0.5\linewidth-0.85em-0.5\arrayrulewidth\relax}@{\hspace{0.5em}}}
  \multicolumn{2}{@{}>{\hspace{0.5em}\raggedright\arraybackslash}p{\dimexpr\linewidth-1em\relax}@{}}{%
    \vspace{0.15em}
    \algheading{Gen}{}
    \algline{01}{$(f,g,F,G,h) \sample \operatorname{TpdGen}(\alpha,q)$}
    \algline{02}{$B \coloneqq \begin{bmatrix} g & -f \\ G & -F \end{bmatrix} \in \mathbb{F}_q^{2n \times 2n}$}
    \algline{03}{return $(sk \coloneqq B,\; pk \coloneqq h)$}
    \vspace{0.15em}
  } \\
  \hline
  \vspace{0.2em}\algscheme{CoreFalcon\textsuperscript{+}}
  \algheading{Sgn$^{+}$}{$(sk=B,m)$}
  \algline{01}{repeat}
  \algindentline{02}{$r\sample\{0,1\}^{k}$}
  \algindentline{03}{$c\coloneqq H(pk,r,m)\in R_{q}$}
  \algindentline{04}{$(s_{1},s_{2}) \sample \operatorname{PreSmp}\!\left(B,s,(c,0)\right)$}
  \algline{05}{until $\lVert(s_{1},s_{2})\rVert_{2}\leq\beta$}
  \algline{06}{$\sigma\coloneqq(r,s_{2}) \in\{0,1\}^{k}\times R$}
  \algline{07}{return $\sigma$}
  \alggap
  \algheading{Ver}{$(pk=h,m,\sigma=(r,s_{2}))$}
  \algline{08}{$c\coloneqq H(pk,r,m)$}
  \algline{09}{$s_{1}\coloneqq c-s_{2}\cdot h\pmod q$}
  \algline{10}{return $\llbracket \lVert(s_{1},s_{2})\rVert_{2}\leq\beta \rrbracket$}
  &
  \vspace{0.2em}\algscheme{KoalaFalcon}
  \algheading{KoalaSgn}{$(sk=B,m)$}
  \algline{01}{repeat}
  \algindentline{02}{$r\sample\{0,1\}^{k}$}
  \algindentline{03}{$c\coloneqq \operatorname{Poseidon}(pk,r,m)\in R_{q}$}
  \algindentline{04}{$(s_{1},s_{2}) \sample \operatorname{PreSmp}\!\left(B,s,(c,0)\right)$}
  \algline{05}{until $\lVert(s_{1},s_{2})\rVert_{2}\leq\beta$}
  \algline{06}{$\sigma\coloneqq(r,s_{2}) \in\{0,1\}^{k}\times R$}
  \algline{07}{return $\sigma$}
  \alggap
  \algheading{KoalaVer}{$(pk=h,m,\sigma=(r,s_{2}))$}
  \algline{08}{$c\coloneqq \operatorname{Poseidon}(pk,r,m)$}
  \algline{09}{$s_{1}\coloneqq c-s_{2}\cdot h\pmod q$}
  \algline{10}{return $\llbracket \lVert(s_{1},s_{2})\rVert_{2}\leq\beta \rrbracket$}
  \\
  \end{tabular}
  \vspace{0.15em}
  \end{minipage}%
  }

  \caption{KoalaFalcon and CoreFalcon\textsuperscript{+} key generation, signing, and verification.
  $\textsc{KoalaFalcon}=(\mathsf{Gen},\mathsf{KoalaSgn},\mathsf{KoalaVer})$ and
  $\CoreFalcon^{+}=(\mathsf{Gen},\mathsf{Sgn}^{+},\mathsf{Ver})$.}
  \label{fig:koalafalcon-construction}
\end{figure}
  
\begin{table}[!htbp]
  \centering
  \setlength{\fboxsep}{9pt}
  \fbox{%
  \begin{minipage}{0.94\textwidth}
  \vspace{0.2em}
  \centering
  \footnotesize
  \renewcommand{\arraystretch}{1.15}
  \begin{tabular}{@{}>{\raggedright\arraybackslash}p{0.46\linewidth}|>{\raggedright\arraybackslash}p{0.46\linewidth}@{}}
  \multicolumn{1}{@{}c|}{\textsc{KoalaFalcon-512}} &
  \multicolumn{1}{c}{\textsc{KoalaFalcon-1024}} \\
  \hline
  $\lambda = 128$ &
  $\lambda = 256$ \\
  $n = 512$ &
  $n = 1024$ \\
  $q = 2^{31} - 2^{24} + 1$ &
  $q = 2^{31} - 2^{24} + 1$ \\
  $Q_s = 2^{64}$ &
  $Q_s = 2^{64}$ \\
  $C_s = 2^{64} + 2^{50}$ &
  $C_s = 2^{64} + 2^{50}$ \\
  $Q_H = 2^{96}$ &
  $Q_H = 2^{96}$ \\
  $\alpha = 1.17$ &
  $\alpha = 1.17$ \\
  $k = 416$ &
  $k = 416$ \\
  $\tau = 1.1$ &
  $\tau = 1.1$ \\
  $\epsilon = 1/\sqrt{Q_s \cdot \lambda}$ &
  $\epsilon = 1/\sqrt{Q_s \cdot \lambda}$ \\
  $\beta = \tau s \sqrt{2n}$ &
  $\beta = \tau s \sqrt{2n}$ \\
  \paramSformula &
  \paramSformula \\
  \end{tabular}
  \vspace{0.2em}
  \end{minipage}%
  }
  \caption{Concrete parameters for KoalaFalcon-512 and KoalaFalcon-1024.}
  \label{tab:koalafalcon-parameters}
\end{table}

\FloatBarrier
\subsection{Security Reduction}\label{modulus-invariance-in-the-reduction}

In this section we discuss implications to the CoreFalcon+ security analysis from
setting modulus $q$ to the KoalaBear prime. 

\paragraph{Requirements on $q$ for the Algebraic Feasibility of the Construction}

First, we discuss whether the KoalaFalcon construction in \hyperref[koalafalcon-construction]{Figure 1} 
is algebraically feasible to instantiate with the KoalaBear prime \(q = 2^{31} - 2^{24} + 1\). This involves ensuring 
the necessary ring-algebraic requirements are in place to support efficient ring multiplication 
and that the desired NTRU lattice structure exists to support trapdoor sampling. Fast ring 
multiplication in Falcon uses the NTT, requiring \(q\) to permit a primitive \(2n\)-th root of 
unity in \(\mathbb{Z}_q\). This is to enable \(X^n + 1\) to split fully into \(n\) linear factors 
over \(\mathbb{Z}_q\). As well, the public and private bases output from trapdoor sampling 
\(B_h, B_{f,g}\) must remain bases of the same NTRU lattice. The properties hold \(q = 2^{31} - 2^{24} + 1\) 
and are proven in \hyperref[appendix-alg-feas]{Appendix~\ref*{appendix-alg-feas}}.

\paragraph{Requirements on $q$ in the $t\text{-}\mathcal{R}\text{-}\mathbf{ISIS}$ Game}

Next, we discuss whether the CoreFalcon+ security reduction requires any properties of \(q\) beyond 
the desired ring setting. At this stage, we are not concerned with hardness estimates, only 
whether the hardness estimates of $\operatorname{Adv}_{q,\alpha,B,\mathcal{A}}^{t\text{-}\mathcal{R}\text{-}\mathbf{ISIS}}$ 
(the advantage of an adversary in solving the $t\text{-}\mathcal{R}\text{-}\mathbf{ISIS}$ game) are
still applicable without modification to the reduction for $q=2^{31} - 2^{24} + 1$.

\[
\operatorname{Adv}_{q,\alpha,B,\mathcal{A}}^{t\text{-}\mathcal{R}\text{-}\mathbf{ISIS}}
:=
\Pr\!\bigl[t\text{-}\mathcal{R}\text{-}\mathbf{ISIS}_{q,\alpha,B}(\mathcal{A}) \Rightarrow 1\bigr].
\]

The CoreFalcon+ security reduction demonstrates adversaries capable of forgery can be reduced to
solvers of the \(t\)-R-ISIS problem depicted in \hyperref[fig:t-r-isis-game]{Figure 2}. The reduction is grounded in the one-wayness of the 
public-to-private key map and standard GPV proof strategies. Rather than arguing the 
applicability of the reduction to concrete instantiations via collision resistance and 
statistical distance of 
public keys to uniformity (as is the case for NTRUSign), CoreFalcon+ uses Rényi divergence 
to quantify the bit loss in security from the real signing and random oracle instantiation 
outputs to the ideal simulated setting. It is assumed \(t\)-R-ISIS is as hard as standard 
SIS and hardness estimates are informed from best-known SIS attacks via the
core-SVP methodology. A hybrid argument upper bounds a more directly applicable NTRU-ISIS 
advantage into an R-ISIS hardness plus decisional-NTRU advantage, arguing that 
\(t\)-R-ISIS is sufficient for the concrete security bounds on unforgeability. Rényi 
divergence is introduced into the bound as a multiplicative factor in the
computational term to quantify the bit loss in security from simulated to real oracles.

\begin{figure}[ht]
  \centering
  \setlength{\fboxsep}{7pt}
  \fbox{%
  \begin{varwidth}{0.36\textwidth}
    \vspace{0.1em}
    \noindent{\scriptsize\textbf{Game $t$-$\mathcal{R}$-ISIS$_{q,\alpha,B}(\mathsf{A})$}}\par
    \vspace{0.08em}
    \hrule
    \vspace{0.12em}
    \algline{01}{$(\cdot,\cdot,\cdot,\cdot,h) \sample \operatorname{TpdGen}(\alpha,q)$}
    \algline{02}{\textbf{for} $i \in [t]$}
    \algindentline{03}{$c_i \sample R_q$}
    \algline{04}{$(j,u,v) \sample \mathsf{A}(h,c_1,\ldots,c_t)$}
    \algline{05}{\textbf{return} $\llbracket u + h\cdot v = c_j \land \lVert(u,v)\rVert_2 \le B \rrbracket$}
    \vspace{0.1em}
  \end{varwidth}%
  }
  \caption{$t$-$\mathcal{R}$-ISIS$_{q,\alpha,B}$ Game.}
  \label{fig:t-r-isis-game}
\end{figure}

The $t\text{-}\mathcal{R}\text{-}\mathbf{ISIS}$ game itself does not impose any strict requirement 
on $q$ aside from the ring $R_q$ existing. Thus, in the next section we turn our attention to how 
this ideal simulation in the ROM is bridged to a real instantiation, tracking exactly where the 
modulus enters into the security bound, and what requirements it must satisfy, term by term. 

\paragraph{Requirements on $q$ for the Security Reduction to $t\text{-}\mathcal{R}\text{-}\mathbf{ISIS}$}\label{sqrtq-scaling-argument}

The reduction from CoreFalcon+ forgery to a $t\text{-}\mathcal{R}\text{-}\mathbf{ISIS}$ introduces 
Rényi-loss terms to quantify distributional differences between the idealized security game and 
a real instantiation. We argue that once the desired ring setting is established, 
the only $q$ dependent property needed to maintain applicability of the CoreFalcon+ security bound 
under different moduli is that the signature standard deviation must be greater than or equal to 
the scaled smoothing parameter of the underlying lattice. While there are $q$ dependent properties,
the CoreFalcon+ formulaic parameters for $\ell_2$ norm threshold $\beta$, target Gram-Schmidt norm 
of the private basis $||B_{f,g}||_{GS},$ and signature standard deviation $s$ are set in such a way 
that this condition is fulfilled for generic $q$. 

\begin{lemma}\label{lem:sqrtq-scaling}
  Fix $n$, $\epsilon$, $\alpha$, $\tau$. Let $\mathrm{TpdGen}(\alpha,q)$ return trapdoor basis $B_{f, g}$ with 
  $\|B_{f,g}\|_{\mathrm{GS}}\le\alpha\sqrt{q}$. Where $s, \beta$ are set as
  \[
    s = \eta_\epsilon(\mathbb{Z}^{2n}) \cdot \alpha \sqrt{q},
    \qquad \,\,
    \eta_\epsilon(\mathbb{Z}^{2n}) = \frac{1}{\pi} \sqrt{\frac{\ln\!\left(4n\left(1+\frac{1}{\epsilon}\right)\right)}{2}},
    \qquad \,\,
    \beta = \tau s \sqrt{2n}
  \]
  Whether the smoothing condition $s\ge\eta_\epsilon(\mathbb{Z}^{2n})\cdot\|B_{f,g}\|_{\mathrm{GS}}$ 
  (and hence $s\ge\eta_\epsilon(\Lambda_{h,q})$) holds is independent of $q$ for every modulus $q$ in 
  which $R_q$ is defined.
\end{lemma}

\begin{proof} 

Notice that \(s \propto \sqrt{q}\) and \(\beta \propto \sqrt{q}\), for fixed \(n, \epsilon, \alpha, \tau\). 
In particular, when normalized by a factor of $\frac{1}{\sqrt{q}}$, the value of 
$\frac{s}{\sqrt{q}} , \frac{\beta}{\sqrt{q}}$ are both independent of \(q\). 


\[
\frac{s}{\sqrt{q}} = \eta_\epsilon(\mathbb{Z}^{2n}) \cdot \alpha
\]

\[
\frac{\beta}{\sqrt{q}} =  \frac{\tau s \sqrt{2n}}{\sqrt{q}}  = \tau \sqrt{2n} \cdot \eta_\epsilon(\mathbb{Z}^{2n}) \cdot \alpha
\]

The trapdoor sampler \(\mathrm{TpdGen}\) determines 
ring elements $f, g, F, G$ defining $B_{f, g}$ satisfying

\[
  B_{f, g} =
  \begin{bmatrix}
  g & -f \\
  G & -F
  \end{bmatrix},
\qquad
||B_{f,g}||_{GS} \le \alpha \sqrt{q}.
\]

The intended smoothing condition on $s$ is $s \ge \eta_\epsilon(\mathbb{Z}^{2n}) \cdot ||B||_{GS}.$
Since both \(s\) and the upper bound on
\(\eta_\epsilon(\mathbb{Z}^{2n}) \cdot ||B||_{GS}\) scale proportional to
\(\sqrt{q}\), changing \(q\) preserves the inequality as long as the trapdoor-quality guarantee is maintained. 

\[
\eta_\epsilon(\mathbb{Z}^{2n})
    \|B_{f,g}\|_{\mathrm{GS}}
\le
\eta_\epsilon(\mathbb{Z}^{2n})
    \alpha\sqrt{q}
=
s.
\]

Thus, for fixed \((n, \epsilon, \alpha)\), satisfaction of the smoothing condition is independent of \(q\),
provided the trapdoor-quality bound
\[
\|B_{f,g}\|_{\mathrm{GS}}\le\alpha\sqrt{q}
\]
continues to hold.

\end{proof}


\begin{theorem}\label{thm:sqrtq-scaling}
  Under the $q$-dependence audit of the CoreFalcon+ reduction given in
  \hyperref[appendix-detailed-loss-term-derivations]{Appendix~\ref*{appendix-detailed-loss-term-derivations}},
  changing the modulus preserves the conditions required for the R\'enyi-loss analysis,
  provided the required cyclotomic ring setting and trapdoor-quality/smoothing conditions remain satisfied.
\end{theorem}

\begin{proof}

A term-by-term audit of the CoreFalcon+ security reduction and security bound 
(in \hyperref[appendix-detailed-loss-term-derivations]{Appendix~\ref*{appendix-detailed-loss-term-derivations}})
finds the only $q$ dependent properties required for the reduction are the cyclotomic ring setting $R_q$
and $s \ge \eta_\epsilon(\mathbb{Z}^{2n}) \cdot ||B||_{GS}$. We have already established in 
\hyperref[lem:sqrtq-scaling]{Lemma ~\ref*{lem:sqrtq-scaling}} that the 
CoreFalcon+ parameterization results in $s \ge \eta_\epsilon(\mathbb{Z}^{2n}) \cdot ||B||_{GS}$ being 
independent of $q$. 
Therefore, so long as $q$ permits the cyclotomic ring structure $R_q$, modification of $q$ does not 
affect the applicability of the reduction.

\end{proof}


\begin{corollary}[Modulus-independent normalized SIS target]\label{thm:sis-target-invariance}
For $m = 2n$,
\[
\frac{\beta}{(\det\Lambda_q^\perp(A))^{1/m}}
=
\frac{\beta}{q^{n/m}}
=
\frac{\beta}{\sqrt{q}},
\]
and therefore the target norm remains constant relative to the determinant scale of the $q$-ary lattice under a change of modulus.
\end{corollary}

\begin{proof}

The same argument holds for \(\beta = \tau s \sqrt{2n}\) with respect to $s$. 
Notice 
\(\beta/s = \tau\sqrt{2n}\) is also independent of \(q\), obeying the same $\sqrt{q}$ scaling behavior. 
Increasing \(q\) raises both \(s\) and \(\beta\) proportionally, so \(\beta\) itself scales as \(\sqrt{q}\),
while the normalized quantity \(\beta / \sqrt{q}\) is independent of \(q\).
The hardness of $t\text{-}\mathcal{R}\text{-}\mathbf{ISIS}$ will depend on 
the shortness of $\ell_2$ norm threshold $\beta$ relative to the natural scale of the lattice.
With \(m = 2n\) the natural scale of a \(q\)-ary lattice is \((\det \Lambda_q^\perp(A))^{1/m} = q^{n/m} = q^{1/2}\). 
So

\[
\frac{\beta}{(\det\Lambda_q^\perp(A))^{1/m}}
=
\frac{\beta}{q^{n/m}}
=
\frac{\beta}{\sqrt{q}}
\]

is the same \(q\)-independent quantity as before, contextualizing why the normalized
SIS target length is stable across the modulus change. The normalized $\beta / \sqrt{q}$ 
remains constant under variation of $q$.

\end{proof}


\paragraph{Lattice Reduction Attacks on q-ary Lattices}\label{lattice-reduction-attacks-on-q-ary-lattices}

Next, to support the findings from the lattice-estimator output, we argue that these parameters result 
in predictable behavior for the estimated hardness of $t\text{-}\mathcal{R}\text{-}\mathbf{ISIS}$
by analyzing lattice reduction attacks on \(q\)-ary lattices. 

Let \(\Lambda_q^\perp(A)\) be an \(m\)-dimensional
\(q\)-ary lattice, where \(q\) is prime, \(A \in \mathbb{Z}_q^{n \times m}\) is a random matrix, 
with linearly independent rows and \(m \ge n\). 
The number of elements in \(\Lambda_q^\perp(A) \mod q\) is \(q^{m-n}\).
The determinant is \(\det(\Lambda_q^\perp(A)) = q^n\).

\[
\Lambda_q^\perp(A) = \{y \in \mathbb{Z}^m : Ay = 0 \mod q\}
\]


By Definition~2.11 of \cite{Menezes2026GentleIntro} 
The volume of \(\Lambda_q^\perp(A)\) is \(\det(\Lambda_q^\perp(A)) = q^n\), and 
corresponds to the sparsity of lattice points in \(\mathbb{R}^m\).
By (Section~3, ~\cite{Regev2009LatticeSurvey}), a heuristic estimate can be made for 
the first successive minimum of \(\Lambda_q^\perp(A)\) as the smallest radius of
an \(m\) dimensional ball with volume \(q^n\)

\[
\lambda_1(\Lambda_q^\perp(A)) \approx q^{n/m} \sqrt{\frac{m}{2\pi e}}
\]

An initial observation is that the determinant of a \(q\)-ary SIS
lattice scales exactly as \(q^n\), and the heuristic estimate from
\cite{Regev2009LatticeSurvey}
of the first successive minimum \(\lambda_1\) scales proportional to
\(q^{n/m}\). For \(m=2n\), the estimated shortest vector in the lattice
\(\Lambda_h\) is approximately

\[
\lambda_1(\Lambda_q^\perp(A)) \approx q^{n/m} \sqrt{\frac{m}{2\pi e}} 
= q^{n/2n} \sqrt{\frac{2n}{2\pi e}}  
= q^{1/2} \sqrt{\frac{n}{\pi e}}  
= \sqrt{\frac{qn}{\pi e}}  
\]

At the time of equation~(1), Section~3 of
\cite{Regev2009LatticeSurvey} (2008), the best known lattice reduction algorithms for determining
short vectors in a random \(m\) dimensional \(q\)-ary lattice
\(\Lambda^\perp_q(A)\) were capable of solving for lengths close to

\[
\min\{q, (\det(\Lambda^\perp_q(A)))^{1/m} \cdot \delta^m\} = \min\{q, q^{n/m}\delta^m\}
\]

Where \(\delta\) depends on the algorithm and is observed by
\cite{GamaNguyen2008} to be typically \(\delta \approx 1.013\) for fast
algorithms, and as low as \(\delta \approx 1.011\) for slower more
precise algorithms. For our parameter choices

\[
\min\{q, q^{n/m}\delta^m\} = \min\{q, q^{n/{2n}}\delta^{2n}\} = \min\{q, \sqrt{q} \cdot \delta^{2n}\} 
\]

Using the script \texttt{bkz\_length.py} in the supporting code, \(\min\{q, q^{n/m}\delta^m\}\) 
is used to compare the signature norm bound \(\beta\) against these estimates.
For the \(\lambda=128, n=512\) setting \(\delta\) would need to be at
most \(1.00387782756709\). For the \(\lambda=256, n=1024\) setting
\(\delta\) would need to be at most \(1.00211437276073\). Meaning for
these style lattice reduction attacks to be effective, \(\delta\) needs
to be below these values of \(\delta\) for the respective parameter
sets. Which is out of practical reach by the analysis of \cite{GamaNguyen2008}.



\subsection{Security Estimates}\label{concrete-parameter-choices}

Now that we have established the viability of the CoreFalcon+ reduction for our 
use case, we re-estimate an adversary's advantage
$\operatorname{Adv}^{(Q_H+1)\text{-R-ISIS}}_{q,\alpha,\beta,B}$
in solving the $(Q_H+1)$-R-ISIS problem under the new modulus.
The CoreFalcon+ approach assumes that $t$-R-ISIS is as hard as SIS
and estimates the resulting computational security using the
\texttt{lattice-estimator} library~\cite{LatticeEstimator}.
We briefly review this methodology to clarify both the origin of the
computational security term and the role played by the reduction to
$t$-R-ISIS before presenting the concrete estimates for our parameter
sets.
It is explicitly assumed that $t$-R-ISIS is as hard as SIS, in part because best known attacks do
not rely on the ring structure of R-ISIS. Estimates of SIS are made using the
core-SVP methodology from section~6 of \cite{Albrecht2015LWE},
which considers dual and primal BKZ / sieve style attacks for predicting the
concrete hardness of SVP style problems. They do not consider enumeration style
attacks because they run in super-exponential time, whereas BKZ runs in exponential time.
This is also the methodology behind the \texttt{SIS.estimate.rough()} function in
\texttt{lattice-estimator}. The core-SVP hardness estimates are:

\begin{center}
\begin{tabular}{@{}rrrrrr@{}}
\toprule
$n$ & $\lambda$ & $q$ & $s$ & $\beta$ & SIS estimate \\
\midrule
512 & 128 & 2130706433 & 69011.59 & $2.4292\cdot 10^6$ & $2^{121.18}$ \\
1024 & 256 & 2130706433 & 70115.84 & $3.4904\cdot 10^6$ & $2^{279.15}$ \\
\bottomrule
\end{tabular}
\end{center}

Running \texttt{scripts/corefalcon.sage} from the supporting code evaluates the Theorem 1 UF-CMA 
bound from a Closer Look at Falcon for the KoalaBear instantiation. 

\begin{center}
\begin{tabular}{@{}rrrr@{}}
\toprule
$n$ & $\lambda$ & SIS bits & UF-CMA bits \\
\midrule
512 & 128 & 121.18 & 117.98 \\
1024 & 256 & 279.15 & 256.00 \\
\bottomrule
\end{tabular}
\end{center}

Rényi and sampler term decomposition from the same run:

\begin{center}
\begin{tabular}{@{}rrrrrrr@{}}
\toprule
$n$ & $\lambda$ & $a_p$ & $a_u$ & $r_p$ & $r_u$ & $p_{\mathrm{PreSmp},\beta}$ \\
\midrule
512 & 128 & 75.35 & 75.35 & 1.00 & 1.00 & 0.999951 \\
1024 & 256 & 150.70 & 150.70 & 1.00 & 1.00 & 0.999999998 \\
\bottomrule
\end{tabular}
\end{center}

% TODO do this if time allows

\subsection{Performance}

Benchmarks for KoalaFalcon-512 signatures were collected on a MacBook Pro with the M4 Max chip 
with 48GB RAM. Throughput and latency for a run of $2^{17}$ signatures. Signatures are $2.1$ KB. 
Multi-threaded benchmarks give each thread a single-threaded signing job, which independently 
run across 16 threads.

\begin{table}[htbp]
  \centering
  \resizebox{0.85\textwidth}{!}{%
  \begin{tabular}{lrrrrrrr}
  \toprule
  Mode
  & Threads
  & $\log_2 N$
  & Keygen (s)
  & Sign total (s)
  & Sign/s
  & Verify total (s)
  & Verify/s \\
  \midrule
  Single-threaded
  & 1
  & 17
  & 0.104
  & 73.286
  & 1788
  & 9.566
  & 13701 \\
  Multi-threaded
  & 16
  & 17
  & 0.092
  & 6.363
  & 20599
  & 0.816
  & 160625 \\
  \bottomrule
  \end{tabular}%
  }
  \caption{Falcon-512 performance, 131,072 signatures.}
  \label{tab:falcon512-signature-performance}
\end{table}

\begin{table}[htbp]
  \centering
  \resizebox{0.85\textwidth}{!}{%
  \begin{tabular}{lrrrrrrr}
  \toprule
  Mode
  & Threads
  & $\log_2 N$
  & Keygen (s)
  & Sign total (s)
  & Sign/s
  & Verify total (s)
  & Verify/s \\
  \midrule
  Single-threaded
  & 1
  & 17
  & 0.795
  & 166.747
  & 786
  & 36.205
  & 3620 \\
  Multi-threaded
  & 16
  & 17
  & 0.388
  & 15.332
  & 8549
  & 3.421
  & 38310 \\
  \bottomrule
  \end{tabular}%
  }
  \caption{Falcon-1024 performance, 131,072 signatures.}
  \label{tab:falcon1024-signature-performance}
\end{table}

Runtime performance and signature size is strictly worse than the original Falcon scheme, 
but this is to be expected given the increase in field size (and therefore the data size 
of each ring coefficient).

\begin{table}[!htbp]
  \centering
  \resizebox{0.85\textwidth}{!}{%
  \begin{tabular}{lrrrrrrr}
  \toprule
  Mode
  & Threads
  & $\log_2 N$
  & Keygen (s)
  & Sign total (s)
  & Sign/s
  & Verify total (s)
  & Verify/s \\
  \midrule
  Single-threaded
  & 1
  & 17
  & 0.010
  & 13.601
  & 9637
  & 1.507
  & 86999 \\
  Multi-threaded
  & 16
  & 17
  & 0.002
  & 1.162
  & 112799
  & 0.129
  & 1017779 \\
  \bottomrule
  \end{tabular}%
  }
  \caption{FN-DSA / Falcon-512 performance, 131,072 signatures.}
  \label{tab:fndsa-falcon512-signature-performance}
\end{table}

\begin{table}[!htbp]
  \centering
  \resizebox{0.85\textwidth}{!}{%
  \begin{tabular}{lrrrrrrr}
  \toprule
  Mode
  & Threads
  & $\log_2 N$
  & Keygen (s)
  & Sign total (s)
  & Sign/s
  & Verify total (s)
  & Verify/s \\
  \midrule
  Single-threaded
  & 1
  & 17
  & 0.011
  & 27.025
  & 4850
  & 2.980
  & 43987 \\
  Multi-threaded
  & 16
  & 17
  & 0.013
  & 2.140
  & 61255
  & 0.247
  & 531447 \\
  \bottomrule
  \end{tabular}%
  }
  \caption{FN-DSA / Falcon-1024 performance, 131,072 signatures.}
  \label{tab:fndsa-falcon1024-signature-performance}
\end{table}

\FloatBarrier